% Successor to Article V: focused F051 incorporation.
% Bilingual fragment paired with the homologous ES file; do not compile in isolation.
\subsection{Prime-vacancy observer and modal resolution}
\label{sec:f051-prime-vacancy-observer}

In the causal chain of the already generated HMT state, the nonadic
product, its reduced multiplicative coproduct, and the spectral projection
\(P_{\rm irr}\) produce the irreducible forms; the evaluation
\(\operatorname{ev}_9\) subsequently recognizes them as ordinary primes.
The \(729/1000\) boundary produces the vacancy calendar. The observer
\(\mathcal C_{\varepsilon,X}\) composes both internal objects with the
logarithmic clock. A conventional sieve enters only as an external reader
for numerical reproduction.

\subsubsection{Calendar and equivalence of the two representations}

Let
\begin{equation}
 \varepsilon_{\rm vac}
 =1-\log_{10}9
 =\log_{10}\!\left(\frac{10}{9}\right),
 \qquad
 z_n=\left\lceil(n+1)\varepsilon_{\rm vac}\right\rceil
      -\left\lceil n\varepsilon_{\rm vac}\right\rceil ,
 \quad n\geq1.
 \label{eq:f051-epsilon-calendar}
\end{equation}
The number \(\varepsilon_{\rm vac}\) is irrational. Indeed, if
\(\varepsilon_{\rm vac}=a/b\in\mathbb Q\), then
\((10/9)^b=10^a\), contradicting unique prime factorization because of the
nonzero power of \(3\). Consequently, neither
\(n\varepsilon_{\rm vac}\) nor
\((n+1)\varepsilon_{\rm vac}\) is an integer for \(n\geq1\), and
\begin{equation}
 z_n=\left\lfloor(n+1)\varepsilon_{\rm vac}\right\rfloor
      -\left\lfloor n\varepsilon_{\rm vac}\right\rfloor
 \in\{0,1\}.
 \label{eq:f051-ceil-floor}
\end{equation}
The restriction \(n\geq1\) is necessary for this equivalence: at
\(n=0\), the ceiling and floor conventions do not agree.

\subsubsection{Typed chain and correlation observer}

For \(X\geq2\), the causal chain can be written without loss of type as
\begin{equation}
 \operatorname{Ran}(P_{\rm irr})
 \xrightarrow{\ \operatorname{ev}_9\ }
 \mathbb P
 \xrightarrow{\ \iota\ }
 \{0,1\}\times\mathbb R_{>0},
 \qquad
 \bigl((z_p,\log p)\bigr)_{p\in\mathbb P_{\leq X}}
 \xrightarrow{\ \mathcal C_{\varepsilon,X}\ }
 \mathbb R,
 \label{eq:f051-typed-chain}
\end{equation}
where
\[
 \iota(p)=(z_p,\log p),
 \qquad
 \mathcal C_{\varepsilon,X}\colon
 \prod_{p\in\mathbb P_{\leq X}}
       (\{0,1\}\times\mathbb R_{>0})\longrightarrow\mathbb R,
 \qquad
 \mathcal C_{\varepsilon,X}\bigl((a_p,\ell_p)_{p\leq X}\bigr)
 =\sum_{p\leq X}(a_p-\varepsilon_{\rm vac})\ell_p .
\]
It follows that
\begin{equation}
 \theta_{\rm vac}(X)=\sum_{p\leq X}z_p\log p,
 \qquad
 \theta(X)=\sum_{p\leq X}\log p,
 \qquad
 R_{\rm vac}(X)=\theta_{\rm vac}(X)
                   -\varepsilon_{\rm vac}\theta(X).
 \label{eq:f051-rvac-chebyshev}
\end{equation}
The exact centering of the observer is
\(\varepsilon_{\rm vac}\theta(X)\); the expression
\(\varepsilon_{\rm vac}X\) belongs to a different normalization. Two
prefixes containing the same number of vacancies can yield different
values of \(R_{\rm vac}\), because the observer retains both prime
position and logarithmic weight.

\subsubsection{Dirichlet readers and analytic domains}

The finite observer and the two Dirichlet readers constitute distinct,
compatible realizations:
\begin{align}
 Z^{+}_{\rm vac}(s)
   &=\sum_{n\geq1}\frac{z_n}{n^s}
     =\varepsilon_{\rm vac}\zeta(s)+H_{\rm vac}(s),
 & H_{\rm vac}&\in\mathcal O(\operatorname{Re}s>0),
 \nonumber\\
 Z^{\times}_{\rm vac}(s)
   &=\prod_p(1-p^{-s})^{-z_p}
     =\zeta(s)^{\varepsilon_{\rm vac}}G_{\rm vac}(s),
 & \operatorname{Re}s&>1,
 \label{eq:f051-dirichlet-readers}\\
 \log G_{\rm vac}(s)
   &=\sum_p(z_p-\varepsilon_{\rm vac})
       \log(1-p^{-s})^{-1},
 & \operatorname{Re}s&>1.
 \nonumber
\end{align}
The first series and the Euler product converge absolutely for
\(\operatorname{Re}s>1\). The partial discrepancy of the mechanical word
is bounded, and Abel summation extends \(H_{\rm vac}\) holomorphically to
\(\operatorname{Re}s>0\). The power
\(\zeta(s)^{\varepsilon_{\rm vac}}\) is defined here through the Euler
logarithm in the half-plane of absolute convergence.

\subsubsection{Fourier resolution and reality of the observer}

For \(k\in\mathbb Z\setminus\{0\}\), define
\begin{equation}
 c_k=\frac{e^{2\pi i k\varepsilon_{\rm vac}}-1}{2\pi i k},
 \qquad
 S_k(X)=\sum_{p\leq X}(\log p)
             e^{2\pi i k p\varepsilon_{\rm vac}}.
 \label{eq:f051-fourier-data}
\end{equation}

\begin{proposition}[Modal decomposition of the prime correlation]
\label{prop:f051-fourier-resolution}
For every \(X\geq2\),
\begin{equation}
 R_{\rm vac}(X)
 =\lim_{K\to\infty}
   \sum_{0<|k|\leq K}c_kS_k(X)
 =\lim_{K\to\infty}
   \sum_{0<|k|\leq K}
   \left(1-\frac{|k|}{K+1}\right)c_kS_k(X).
 \label{eq:f051-fourier-pair}
\end{equation}
The first limit is defined through symmetric partial sums and the second
through Fej\'er means; these are exactly the two asserted notions of
convergence. Absolute convergence in \(k\) lies outside both the
hypotheses and the conclusion.
\end{proposition}

\begin{proof}
Let
\[
 f_{\varepsilon}(x)
 =\mathbf1_{(1-\varepsilon_{\rm vac},1)}(\{x\})
  -\varepsilon_{\rm vac}.
\]
Its nonzero Fourier coefficients are precisely \(c_k\). The irrationality
of \(\varepsilon_{\rm vac}\) prevents
\(p\varepsilon_{\rm vac}\) from reaching an endpoint of the interval; by
\eqref{eq:f051-ceil-floor},
\(f_{\varepsilon}(p\varepsilon_{\rm vac})
 =z_p-\varepsilon_{\rm vac}\). Dirichlet's theorem gives convergence of
the symmetric sums at these continuity points, and Fej\'er's theorem gives
convergence of the means. Finiteness of the set \(p\leq X\) permits
interchange of the limit with the weighted prime sum under symmetric and
Fej\'er convergence. The estimate \(|c_k|=O(|k|^{-1})\) does not imply
absolute summability and therefore does not permit an absolute
rearrangement.

Finally,
\(c_{-k}=\overline{c_k}\) and
\(S_{-k}(X)=\overline{S_k(X)}\). The modes \(k\) and \(-k\) contribute
\(2\operatorname{Re}(c_kS_k(X))\), which directly recovers the reality of
\(R_{\rm vac}(X)\).
\end{proof}

\subsubsection{Numerical witnesses and independent reproduction}

Table~\ref{tab:f051-rvac-powers} presents the six available finite
cutoffs. The values of \(R_{\rm vac}\) reproduce every stored digit, and
the normalized column retains the rounding of the numerical source.

\begin{table}[htbp]
\centering
\small
\begin{tabular}{rrrrr}
\toprule
\(X\) & \(\pi(X)\) & \(\pi_{\rm vac}(X)\) & \(R_{\rm vac}(X)\)
 & \(R_{\rm vac}(X)/\sqrt X\)\\
\midrule
\(10^3\) & 168 & 7 & \(-4.553590011317908\) & \(-0.143997159663565\)\\
\(10^4\) & 1\,229 & 61 & \(41.334631595982586\) & \(0.413346315959826\)\\
\(10^5\) & 9\,592 & 427 & \(-141.1579619980301\) & \(-0.446380669781268\)\\
\(10^6\) & 78\,498 & 3\,595 & \(32.21713602110386\) & \(0.032217136021104\)\\
\(10^7\) & 664\,579 & 30\,384 & \(-418.47581825715594\) & \(-0.132333673139529\)\\
\(10^8\) & 5\,761\,455 & 263\,841 & \(4021.0180144347314\) & \(0.402101801443473\)\\
\bottomrule
\end{tabular}
\caption{Prime-vacancy observer at six finite cutoffs.}
\label{tab:f051-rvac-powers}
\end{table}

The complete thirty-mode table is stored in
\texttt{S\_k\_1\_30\_X1e8.csv}. Its five largest values of
\(|S_k(10^8)|\), sorted by modulus, are:

\begin{table}[htbp]
\centering
\scriptsize
\setlength{\tabcolsep}{3.5pt}
\begin{tabular}{rrrrrr}
\toprule
\(k\) & \(\operatorname{Re}S_k\) & \(\operatorname{Im}S_k\)
 & \(|S_k|\) & \(|S_k|/\sqrt X\) & \(|S_k|/(\sqrt X\log X)\)\\
\midrule
19 & 30646.583055876723 & -39786.84285255695 & 50221.56824686792 & 5.0221568246867925 & 0.2726368745267788\\
24 & -1325.0060816425494 & 47631.087438932576 & 47649.51344695589 & 4.764951344695589 & 0.25867400944234675\\
21 & 1214.7582244307505 & -44824.22356680968 & 44840.68081453668 & 4.484068081453668 & 0.24342575303172864\\
26 & -26879.47488364538 & 33364.6769815865 & 42845.161221614406 & 4.284516122161441 & 0.23259271368502904\\
1 & 40391.94841735065 & 13384.21089257329 & 42551.693246765084 & 4.255169324676508 & 0.23099956965887428\\
\bottomrule
\end{tabular}
\caption{Five largest modes in the sample \(1\leq k\leq30\) at
\(X=10^8\). On this sample domain,
\(\lvert S_{19}(10^8)\rvert\) is the maximum.}
\label{tab:f051-top-five-modes}
\end{table}

The independent calculation of the first four rows, obtained from
\eqref{eq:f051-epsilon-calendar} with a sieve used as a verification
reader, reproduces the values through \(10^6\) with absolute errors below
\(2.1\times10^{-11}\). The rows at \(10^7\) and \(10^8\), together with
the thirty modes, are finite numerical witnesses of the historical
calculation. The projection \(P_{\rm irr}\) in
\eqref{eq:f051-typed-chain} retains the role of internal generator.

\subsubsection{Analytic condition for the critical bound}

The estimate
\begin{equation}
 R_{\rm vac}(X)=O\!\left(X^{1/2}\log^B X\right)
 \label{eq:f051-conditional-bound}
\end{equation}
would control the centered prime term of the factor \(G_{\rm vac}\). Its
proof requires estimates for \(S_k(X)\) that are uniform in \(k\) and
compatible with symmetric or Fej\'er summation. The tables characterize a
finite window and do not provide that uniformity. The proved result is the
exact modal identity for each finite \(X\); the critical asymptotic bound
remains conditional on uniform modal estimates.
